\documentclass[11pt,a4paper]{article}
\usepackage[T1]{fontenc}
\usepackage{lmodern,amsmath,amssymb,amsthm,mathtools,booktabs,microtype}
\usepackage[margin=25mm]{geometry}
\usepackage[hidelinks]{hyperref}
\hypersetup{pdftitle={Boolean blocks and the weak Lefschetz property of whiskered and chordal graph algebras},pdfauthor={Anonymous}}
\newtheorem{theorem}{Theorem}[section]
\newtheorem{lemma}[theorem]{Lemma}
\newtheorem{corollary}[theorem]{Corollary}
\newtheorem{proposition}[theorem]{Proposition}
\newtheorem{maintheorem}{Theorem}
\renewcommand{\themaintheorem}{\Alph{maintheorem}}
\theoremstyle{definition}
\newtheorem{definition}[theorem]{Definition}
\newtheorem{example}[theorem]{Example}
\newtheorem{remark}[theorem]{Remark}
\newcommand{\kk}{\mathbb K}
\newcommand{\rk}{\operatorname{rank}}
\newcommand{\cok}{\operatorname{coker}}
\newcommand{\B}{\mathcal B}
\newcommand{\dual}{^{\vee}}
\title{Boolean blocks and the weak Lefschetz property\\ of whiskered and chordal graph algebras}
\author{Anonymous\\[3pt]{\small DOI \href{https://doi.org/10.5281/zenodo.24540630}{10.5281/zenodo.24540630}}}
\date{November 2026}
\begin{document}
\maketitle

\begin{abstract}
Holleben and Nicklasson conjectured that whiskering a graph whose independence number is at least three always produces an Artinian algebra without the weak Lefschetz property. We prove this conjecture over every field, and we determine exactly when the weak and strong Lefschetz properties hold. If $G$ has $n$ vertices and independence number $\alpha$, the whiskered algebra has the weak property exactly when $\alpha\le1$, or $n$ is odd and $\alpha=2$, and the characteristic is zero or exceeds $\lceil n/2\rceil$. It has the strong property exactly when $\alpha\le1$ and the characteristic is zero or exceeds $n$. The proof rests on one change of variables. After it, multiplication by the sum of the variables splits into shifted copies of the Boolean algebra, one for each independent set of $G$, and a known diagonal form of subset-inclusion matrices then gives every rank in every characteristic. The same splitting settles the Lefschetz properties of the Gorenstein algebras that Holleben and Nicklasson built from whiskered graphs, and gives the generic ranks for higher truncation exponents. For chordal graphs, removing a simplicial vertex gives a similar splitting. From it we obtain a criterion for both properties in every characteristic and a uniform, largely computer-free proof of the published classifications for paths and lollipop graphs, now extended to positive characteristic.
\end{abstract}

\section{Introduction}

This section explains the problem, what was known, what this paper proves and how the paper is organised.

\paragraph{The problem.}
Let $A=A_0\oplus A_1\oplus\cdots\oplus A_D$ be a standard graded Artinian algebra over a field $\kk$, so that $A$ is generated by $A_1$ and $A_D\neq0$ is its top degree. A linear map has \emph{maximal rank} if it is injective or surjective. The algebra $A$ has the \emph{weak Lefschetz property} (WLP) if some $\ell\in A_1$ makes every multiplication map $\times\ell:A_i\to A_{i+1}$ have maximal rank, and the \emph{strong Lefschetz property} (SLP) if some $\ell$ makes every map $\times\ell^{\,j-i}:A_i\to A_j$, $i<j$, have maximal rank. These properties abstract the hard Lefschetz theorem of algebraic geometry. They force the Hilbert function $i\mapsto\dim A_i$ to be unimodal, and for algebras built from combinatorial data they encode Sperner-type statements about the data \cite{Stanley}. Deciding which natural families have them is an active programme; see the introductions of \cite{HN,Cooper,NT}.

A rich source of examples comes from graphs. For a finite simple graph $H$ with vertex variables $z_v$, the algebra
\[
 R(H)=\kk[z_v:v\in V(H)]/(z_v^2,\ z_vz_w:\ vw\in E(H))
\]
has a basis of squarefree monomials $z_I=\prod_{v\in I}z_v$ indexed by the independent sets $I$ of $H$, so its Hilbert function is the independence sequence of $H$. Whiskered graphs are the most studied case. \emph{Whiskering} a graph $G$ on vertices $1,\dots,n$ attaches a new leaf $y_i$ to each vertex $x_i$. Writing $x_i,y_i$ also for the variables,
\begin{equation}\label{AG}
 A_G=R(w(G))=\kk[x_1,\dots,x_n,y_1,\dots,y_n]/(x_i^2,\ y_i^2,\ x_iy_i,\ x_ix_j:\ \{i,j\}\in E(G)).
\end{equation}
Whiskered graphs have Cohen--Macaulay edge ideals, their algebras are level (the socle is concentrated in the top degree), and their independence complexes are balls or spheres \cite{HN,Cooper}. The \emph{independence number} $\alpha(G)$ is the largest size of an independent set of $G$.

\paragraph{What was known.}
Cooper, Faridi, Holleben, Nicklasson and Van Tuyl \cite{Cooper} showed that multiplication by a general linear form on $A_G$ has maximal rank in degrees $1$ and $n-1$ when the characteristic is not two, and is injective below degree $n/2$ in characteristic zero, and that some whiskered graphs fail maximal rank in every intermediate degree. Holleben and Nicklasson \cite{HN} proved the SLP when $G$ is complete (and with higher truncation exponents), the WLP in characteristic zero when $\alpha(G)\le2$ and the socle degree is odd, and failure of the WLP when $\alpha(G)\ge n/3+2$, hence, with computer checks for small orders, for all bipartite graphs. They conjectured \cite[Conjecture~1.3 and Conjecture~3.15]{HN} that the WLP fails whenever $\alpha(G)\ge3$, and checked it by computer for graphs on at most seven vertices; an eight-vertex verification in characteristic zero has been posted \cite{Zeng}. They used these algebras to build G-quadratic Gorenstein algebras failing the WLP \cite[Theorem~4.6]{HN}. For general graphs $H$, the WLP of $R(H)$ has been classified in characteristic zero for paths and cycles \cite{NT}, lollipop graphs \cite{TPS} and tadpole graphs \cite{TPS2}, each by an induction specific to the family.

\paragraph{What is new.}
Write $b_n=1$ if $n$ is even and $b_n=2$ if $n$ is odd.

\begin{maintheorem}[Theorem~\ref{base}]\label{thmA}
Let $G$ be a graph on $n$ vertices with independence number $\alpha$, and let $\kk$ have characteristic $p\ge0$. Then $A_G$ has the WLP if and only if $\alpha\le b_n$ and either $p=0$ or $p>\lceil n/2\rceil$. It has the SLP if and only if $\alpha\le1$ and either $p=0$ or $p>n$. In particular the conjecture of Holleben and Nicklasson holds over every field.
\end{maintheorem}

In words: in characteristic zero only complete graphs, and graphs of odd order with independence number two, give the WLP; every non-complete graph of even order fails. The theorem holds for any simplicial complex $\Delta$ in place of the independence complex of $G$ (Section~\ref{sec:split}), and it gives the rank of every power $L^{j-i}:(A_G)_i\to(A_G)_j$ of the sum of the variables $L$ in every characteristic (equation~\eqref{rankformula} and Theorem~\ref{split}).

\begin{maintheorem}[Theorem~\ref{ideal}, Corollary~\ref{tg}]\label{thmB}
The Gorenstein algebra $T$ obtained from $A_G$ by Nagata idealisation, which for non-complete $G$ Holleben and Nicklasson realised in characteristic zero as a G-quadratic Perazzo algebra, has the WLP if and only if $\alpha\le b_n$ and either $p=0$ or $p>\lceil (n+1)/2\rceil$, and the SLP if and only if $\alpha\le1$ and either $p=0$ or $p>n+1$. A single explicit linear form attains every generic rank on $T$ in characteristic zero.
\end{maintheorem}

\begin{maintheorem}[Theorems~\ref{chordal} and \ref{chordalcrit}]\label{thmC}
For a chordal graph $H$, the algebra $R(H)$ splits as a graded module over $\kk[L]$ into shifted Boolean algebras, one for each maximal independent set, and the multiset of these blocks depends only on $H$. The WLP and SLP of $R(H)$ in characteristic zero are decided by the set of block types alone, by an explicit overlap criterion. In characteristic $p>0$ the WLP holds if and only if it holds in characteristic zero and $p>\lceil\alpha(H)/2\rceil$, and the SLP holds if and only if it holds in characteristic zero and $p>\alpha(H)$.
\end{maintheorem}

As applications we reprove the characteristic-zero classification of Nguyen and Tran for paths \cite{NT} and of Tran Quang Hoa, Nguyen Duy Phuoc and Tran Nguyen Thanh Son for lollipop graphs \cite{TPS}, with short explicit failure witnesses in place of family-specific inductions, extend both to every characteristic, and add the SLP (Corollaries~\ref{paths} and~\ref{lollipops}).

\paragraph{The key idea.}
Replace $y_i$ by $u_i=x_i+y_i$. In the new variables the defining ideal of $A_G$ is still monomial, and $L=\sum_i(x_i+y_i)$ becomes $\sum_iu_i$, which never changes the set of $x$-variables present in a monomial. So $L$ acts separately on the span of the monomials $x_Su_R$ with a fixed independent set $S$, and on that span it acts as the sum of the variables of a Boolean algebra on the $n-|S|$ remaining indices. For complete graphs this block structure was already visible, in triangular form, in \cite[Lemma~3.1]{HN}; the change of variables makes it block-diagonal for every graph. The kernel of one block and the cokernel of another then explain failure of maximal rank exactly. For chordal graphs the same substitution at a simplicial vertex, iterated, plays the role of the whiskers; for flag complexes Theorem~\ref{split} is the special case in which every eliminated simplicial vertex is a whisker.

\paragraph{Roadmap.}
Section~\ref{sec:prelim} fixes notation and collects two tools: generic forms over arbitrary fields, and ranks of powers in the Boolean algebra. Section~\ref{sec:split} proves the integral splitting. Section~\ref{sec:wlp} proves Theorem~\ref{thmA} and illustrates it on the published eight-vertex example. Section~\ref{sec:gor} treats the Gorenstein idealisation, Section~\ref{sec:higher} higher truncation exponents, and Section~\ref{sec:chordal} chordal graphs, paths and lollipops. Section~\ref{sec:scope} describes the computational checks and the limits of the method.

\section{Notation and two tools}\label{sec:prelim}

This section fixes notation for simplicial complexes, Boolean algebras and graded Jordan types, then proves that the sum of the variables is a best possible linear form over every field and records the ranks of powers in a Boolean algebra.

\paragraph{Complexes and algebras.}
Throughout, $\kk$ is a field of characteristic $p\ge0$. A simplicial complex $\Delta$ on $[n]=\{1,\dots,n\}$ is a family of subsets (faces) closed under taking subsets, with $\varnothing\in\Delta$. Let $a=\max\{|S|:S\in\Delta\}$ and let $f_s$ be the number of faces of size $s$, so $f_0=1$. For $S\subseteq[n]$ write $x_S=\prod_{i\in S}x_i$. The \emph{whiskered face algebra} of $\Delta$ is
\[
 A_\Delta=\kk[x_1,\dots,x_n,y_1,\dots,y_n]/(x_i^2,\ y_i^2,\ x_iy_i,\ x_F:\ F\notin\Delta).
\]
When $\Delta=\operatorname{Ind}(G)$ is the independence complex of a graph, the nonfaces are generated by edges, $A_\Delta=A_G$ and $a=\alpha(G)$. In the language of \cite{HN}, $A_\Delta$ is the algebra $A(w(\Delta))$ of the whiskered complex. Its monomial basis consists of the $x_Sy_R$ with $S\in\Delta$ and $R\cap S=\varnothing$, so its top degree is $n$. The presentation makes sense over $\mathbb Z$, and we use the same name for the $\mathbb Z$-algebra. Let $L=\sum_i(x_i+y_i)$.

\paragraph{Boolean blocks.}
For a finite set $J$ let $\B_J=\mathbb Z[u_j:j\in J]/(u_j^2:j\in J)$ with the operator $U=\sum_{j\in J}u_j$; we write $\B_m$ when $|J|=m$. Its degree-$r$ part has the basis $u_R$, $|R|=r$, and dimension $\binom mr$. For a graded module $M$, $M\langle s\rangle$ denotes $M$ with its degree $0$ moved to degree $s$. The \emph{twice-centre} of a block $\B_m\langle s\rangle$ is $2s+m$; its Hilbert function is symmetric about half of it. We use $\binom mr=0$ for $r\notin[0,m]$.

\paragraph{Graded Jordan type.}
Over a field, a graded module $M$ over $\kk[L]$ with $L$ of degree one decomposes into \emph{strings}: cyclic summands $\kk[L]v/(L^{k}v)$ spanning one dimension in each of the degrees $\deg v,\dots,\deg v+k-1$. The multiset of degree intervals of the strings is the \emph{graded Jordan type} of $L$; it is equivalent to knowing the ranks of all maps $L^{j-i}:M_i\to M_j$.

\begin{lemma}[Generic forms]\label{generic}
Let $I$ be an ideal of $P=\kk[w_1,\dots,w_N]$ generated by monomials, let $A=P/I$ be Artinian and $L=\sum_kw_k$. Over any field $\kk$, including finite fields, every linear form $h\in A_1$ satisfies
\[
 \rk(h^{j-i}:A_i\to A_j)\le\rk(L^{j-i}:A_i\to A_j)\qquad(0\le i<j).
\]
\end{lemma}
\begin{proof}
Write $h=\sum_kc_kw_k$. Over the rational function field $\kk(t)$ put $c_k'=c_k$ if $c_k\ne0$ and $c_k'=t$ otherwise, and $h_t=\sum_kc'_kw_k$. The diagonal substitution $w_k\mapsto c'_kw_k$ is an automorphism of $P\otimes\kk(t)$ preserving every monomial ideal, so it carries $L$ to $h_t$ and $\rk_{\kk(t)}(h_t^{j-i})=\rk_{\kk(t)}(L^{j-i})=\rk_\kk(L^{j-i})$, the last because $L$ has coefficients in the prime field. The matrix of $h_t^{j-i}$ in the monomial basis has entries in $\kk[t]$ and specialises at $t=0$ to the matrix of $h^{j-i}$; a minor that is nonzero at $t=0$ is a nonzero polynomial, so $\rk(h^{j-i})\le\rk(h_t^{j-i})$.
\end{proof}

For the WLP of monomial algebras this reduction to $L$ is standard; Cooper et al.\ record it over an arbitrary field \cite[Lemma~2.2 and the remark after it]{Cooper}. We include the short proof because it covers all powers of $L$ at once, which the SLP statements and the prime-characteristic thresholds below use. Consequently $A$ has the WLP (SLP) over $\kk$ if and only if $L$ has the corresponding maximal-rank property, and in characteristic zero the graded Jordan type of $L$ is the generic one.

\paragraph{Ranks in a Boolean algebra.}
For $0\le r<t\le m$ the matrix of $U^{t-r}:(\B_m)_r\to(\B_m)_t$ in the monomial bases is $(t-r)!$ times the inclusion matrix $W_{rt}(m)$ of $r$-subsets against $t$-subsets, because $U^{d}u_R=d!\sum u_{R'}$ over the $R'\supseteq R$ with $|R'\setminus R|=d$. A \emph{diagonal form} of an integer matrix $M$ is a diagonal matrix $EMF$ with $E,F$ unimodular. By Wilson's theorem \cite[Theorem~2]{Wilson}, if $r\le\min\{t,m-t\}$, then $W_{rt}(m)$ has a diagonal form with entries $\binom{t-h}{r-h}$ of multiplicity $\binom mh-\binom m{h-1}$, $0\le h\le r$; another proof is given by Ghorbani, Khosrovshahi, Maysoori and Mohammad-Noori \cite[Corollary~3]{GKMM}. Wilson also records the resulting rank of $W_{rt}(m)$ modulo a prime \cite[Theorem~1]{Wilson}. If instead $t>m-r$, complementation of subsets identifies $W_{rt}(m)$ with the transpose of $W_{m-t,m-r}(m)$, to which the same statement applies. Put $d=t-r$ and $b=\min(r,m-t)$. In both cases $\binom{t-h}{r-h}=\binom{b+d-h}{d}$ after the relabelling, so $U^{d}:(\B_m)_r\to(\B_m)_t$ has a diagonal form with nonzero entries
\begin{equation}\label{wil}
 d!\binom{b+d-h}{d}\quad\text{with multiplicity}\quad \binom mh-\binom m{h-1},\qquad 0\le h\le b .
\end{equation}
These multiplicities are positive (since $b\le m/2$) and add up to $\binom mb=\min\{\binom mr,\binom mt\}$. Hence:
\begin{itemize}
\item over a field of characteristic $p$, the rank of $U^d$ counts the entries of \eqref{wil} not divisible by $p$, and $U^d$ has maximal rank if and only if no entry is divisible by $p$;
\item in characteristic zero, or if $p>m$, every power of $U$ has maximal rank: it is injective when $r+t\le m$ and surjective when $r+t\ge m$;
\item for adjacent degrees ($d=1$) the entries are $1,2,\dots,b+1$, and $b+1\le\lceil m/2\rceil$, with equality for $r=\lfloor (m-1)/2\rfloor$.
\end{itemize}
In characteristic zero the Jordan type of $U$ on $\B_m$ is the familiar symmetric one: $\binom mh-\binom m{h-1}$ strings occupying degrees $h,\dots,m-h$ for $0\le h\le\lfloor m/2\rfloor$. Indeed $U^{m-2h}:(\B_m)_h\to(\B_m)_{m-h}$ is bijective, and the strings may be started on a complement of $U((\B_m)_{h-1})$ in $(\B_m)_h$.

\section{The integral splitting}\label{sec:split}

This section proves that, after one change of variables, $L$ acts on $A_\Delta$ as a direct sum of shifted Boolean operators, one for each face of $\Delta$.

\begin{theorem}\label{split}
As graded $\mathbb Z[L]$-modules
\[
 A_\Delta=\bigoplus_{S\in\Delta}x_S\,\B_{[n]\setminus S}
 \;\simeq\;\bigoplus_{s=0}^{a}\B_{n-s}\langle s\rangle^{\oplus f_s},
\]
with $L$ acting on each summand $x_S\B_{[n]\setminus S}$ as $U=\sum_{i\notin S}u_i$, where $u_i=x_i+y_i$. The splitting persists after any change of coefficient ring.
\end{theorem}
\begin{proof}
The substitution $y_i\mapsto u_i-x_i$ and its inverse have integer coefficients, so they give mutually inverse automorphisms of $\mathbb Z[x,y]$. In the new variables the defining ideal is generated by
\[
 x_i^2,\qquad x_iu_i,\qquad u_i^2,\qquad x_F\ (F\notin\Delta),
\]
because $x_iy_i=x_iu_i-x_i^2$ and $y_i^2=u_i^2-2x_iu_i+x_i^2$, while conversely $x_iu_i=x_i^2+x_iy_i$ and $u_i^2=x_i^2+2x_iy_i+y_i^2$. This ideal is monomial, so the monomials $x_Su_R$ with $S\in\Delta$ and $R\subseteq[n]\setminus S$ form a $\mathbb Z$-basis of $A_\Delta$. Since $L=\sum_iu_i$ and $u_ix_Su_R$ is $0$ if $i\in S\cup R$ and $x_Su_{R\cup\{i\}}$ otherwise, multiplication by $L$ preserves the index $S$ and acts on the span of the $x_Su_R$ as $U$ acts on $\B_{[n]\setminus S}$, shifted by $|S|$.
\end{proof}

\begin{remark}
For $\Delta=\operatorname{Ind}(K_n)$ the faces are $\varnothing$ and the singletons, and the splitting is the block-diagonal form of the block-triangular decomposition in \cite[Lemma~3.1]{HN}. The off-diagonal block there disappears in the new variables. The splitting uses that every index $i$ has its own leaf $y_i$: Holleben and Nicklasson's very well-covered example \cite[Example~3.18]{HN}, which has relations $x_iy_j$ with $i\ne j$, has no analogue of the substitution, and it does have the WLP despite an independence number of four on its cover vertices.
\end{remark}

Combining Theorem~\ref{split} with \eqref{wil}, with $m=n-s$, $r=i-s$, $t=j-s$ for each face size $s$, gives a diagonal form over $\mathbb Z$ of every map $L^{j-i}:(A_\Delta)_i\to(A_\Delta)_j$, hence its rank over every field: count the entries not divisible by $p$. In particular the face numbers $f_s$ alone determine all these operators up to change of bases. In characteristic zero, or if $p>n$, every block has maximal rank and
\begin{equation}\label{rankformula}
 \rk(L^{j-i}:A_i\to A_j)=\sum_{s=0}^af_s\min\left\{\binom{n-s}{i-s},\binom{n-s}{j-s}\right\}.
\end{equation}
In these characteristics the graded Jordan type of $L$ (in characteristic zero, by Lemma~\ref{generic}, the Jordan type of a general linear form) consists, for each face of size $s$ and each $0\le h\le\lfloor(n-s)/2\rfloor$, of $\binom{n-s}h-\binom{n-s}{h-1}$ strings in degrees $s+h,\dots,n-h$.

\section{The weak and strong Lefschetz properties}\label{sec:wlp}

This section determines exactly which maps $L^{j-i}$ fail to have maximal rank, and deduces Theorem~\ref{thmA}.

The mechanism is simple. A block $\B_{n-s}\langle s\rangle$ has Hilbert function symmetric about $(n+s)/2$, so in characteristic zero its map from degree $i$ to degree $j$ is injective when $i+j\le n+s$ and surjective when $i+j\ge n+s$. A direct sum of maps has maximal rank only if all summands are injective or all are surjective. Blocks for larger faces sit higher, so a map can be injective on the empty-face block and surjective on a large-face block, or the reverse, and then fail.

\begin{theorem}\label{base}
Let $\kk$ have characteristic $p\ge0$.
\begin{enumerate}
\item If $p=0$ or $p>n$, then for $0\le i<j\le n$ the map $L^{j-i}:(A_\Delta)_i\to(A_\Delta)_j$ fails to have maximal rank if and only if
\begin{equation}\label{strip}
 n<i+j<n+a .
\end{equation}
\item $A_\Delta$ has the WLP if and only if $a\le b_n$ and either $p=0$ or $p>\lceil n/2\rceil$.
\item $A_\Delta$ has the SLP if and only if $a\le1$ and either $p=0$ or $p>n$.
\end{enumerate}
\end{theorem}
\begin{proof}
By Lemma~\ref{generic} it suffices to study $L$: a map $h^{j-i}$ can have maximal rank only if $L^{j-i}$ has, since $\rk h^{j-i}\le\rk L^{j-i}$.

\emph{Step 1: the strip.} Assume $p=0$ or $p>n$, so every block map has maximal rank and its direction is decided by dimensions, as described above (a block map with zero source counts as injective, and one with zero target as surjective). If $i+j\le n$ every block is injective; if $i+j\ge n+a$ every block is surjective. Suppose $n<i+j<n+a$. On the empty-face block $\B_n$ the source $\binom ni=\binom n{n-i}$ is strictly larger than the target $\binom nj$, because $\max(i,n-i)<j\le n$ and binomial coefficients $\binom nk$ decrease strictly for $k\ge n/2$; so this block, hence $L^{j-i}$, has a nonzero kernel. Now put $s=i+j-n+1$. Then $2\le s\le a$, so a face of size $s$ exists by downward closure. Its block has size $m=n-s$, source $\binom{m}{i-s}$ and target $\binom{m}{j-s}$, and $(i-s)+(j-s)=m-1$, so $\binom m{j-s}=\binom m{i-s+1}$ and $0\le j-s\le m$. If $i<s$ the source is zero and the target is not. Otherwise $2(i-s)+1\le(i-s)+(j-s)<m$, so $\binom m{i-s}<\binom m{i-s+1}$. Either way this block has a nonzero cokernel, and $L^{j-i}$ has rank smaller than both $\dim A_i$ and $\dim A_j$.

\emph{Step 2: WLP and SLP in characteristic zero or $p>n$.} For adjacent degrees $j=i+1$ the sum $i+j$ is odd, and an odd integer lies strictly between $n$ and $n+a$ exactly when $a\ge3$, or $a=2$ and $n$ is even, that is, when $a>b_n$; such an $i$ satisfies $i\le n-1$. For the SLP, if $a\ge2$ (so $n\ge2$) the pair $i=1$, $j=n$ lies in the strip; if $a\le1$ the strip contains no integer.

\emph{Step 3: small characteristic, WLP.} For adjacent degrees the entries of \eqref{wil} on a block of size $m\le n$ are at most $\lceil m/2\rceil\le\lceil n/2\rceil$, so if $p>\lceil n/2\rceil$ every adjacent block map has the same rank as in characteristic zero and Step~2 applies. If $2\le p\le\lceil n/2\rceil$, take $i=p-1$, $j=p$. Then $2p-1\le n$, so on the empty-face block $b=p-1$ and \eqref{wil} contains the entry $p$: this block map is not injective although its source is no larger than its target. Since a block map without maximal rank has both a kernel and a cokernel, $L$ fails maximal rank from degree $p-1$ to degree $p$.

\emph{Step 4: small characteristic, SLP.} If $0<p\le n$, then $h^p=\sum c_k^pw_k^p=0$ for every linear form $h$ by the Frobenius identity and $w_k^2=0$, while $(A_\Delta)_p\ne0$; so $h^p:(A_\Delta)_0\to(A_\Delta)_p$ is not injective. If $p>n$ every entry $d!\binom{b+d-h}d$ in \eqref{wil} is a product of integers at most $n$, hence prime to $p$, and Step~2 applies.
\end{proof}

\begin{proof}[Proof of Theorem~\ref{thmA}]
Apply Theorem~\ref{base} to $\Delta=\operatorname{Ind}(G)$, for which $a=\alpha(G)$. When $\alpha(G)\ge3$ we have $a>b_n$, so the WLP fails over every field, which is the conjecture of Holleben and Nicklasson.
\end{proof}

\begin{example}[the published eight-vertex example]\label{ex8}
Let $G$ be $K_8$ with the three edges of a triangle removed, the example of \cite[Example~3.17]{HN}. Then $f=(1,8,3,1)$ and $a=3$, so the strip \eqref{strip} is $8<i+j<11$ and the only adjacent failure is from degree $4$ to degree $5$. Block by block, in characteristic zero:
\[
\begin{array}{c|rrrr|r}
\text{face size }s&0&1&2&3&\text{total}\\ \hline
\dim A_4&70&280&45&5&400\\
\dim A_5&56&280&60&10&406\\
\rk(L:A_4\to A_5)&56&280&45&5&386
\end{array}
\]
The kernel ($400-386=14$) comes from the empty-face block and the cokernel ($406-386=20$) from the three blocks of size two and the block of size three. The dimensions $400$ and $406$ agree with the Hilbert series in \cite{HN}, where this map is reported not injective.
\end{example}

\begin{example}
For $n=3$ and $G$ with a single edge, $f=(1,3,2)$, $a=2=b_3$ and $A_G\simeq\B_3\oplus\B_2\langle1\rangle^{\oplus3}\oplus\B_1\langle2\rangle^{\oplus2}$ has the WLP in characteristic zero and in characteristic $p>2$, but not in characteristic $2$. The SLP fails: $L^2:A_1\to A_3$ has rank $1+3=4$ between two spaces of dimension six, and the two blocks $\B_1\langle2\rangle$ contribute the cokernel.
\end{example}

\section{Gorenstein idealisations}\label{sec:gor}

This section shows that the splitting extends to the Gorenstein algebras built from $A_\Delta$, and proves Theorem~\ref{thmB}.

Write $A=A_\Delta$ over $\kk$ and $D=n+1$. Let $(-)\dual$ denote the dual vector space, extended by zero outside $[0,n]$, and let $A$ act on $A\dual$ by $(b\mu)(w)=\mu(bw)$. The \emph{Nagata idealisation} of $A$ by its graded dual is
\[
 T=\bigoplus_{i=0}^{D}T_i,\qquad T_i=A_i\oplus(A_{D-i})\dual,\qquad (b,\lambda)(c,\mu)=(bc,\ b\mu+c\lambda).
\]
The pairing $T_i\times T_{D-i}\to T_D=(A_0)\dual\cong\kk$ given by multiplication is perfect, so $T$ is Gorenstein with socle degree $D$; this is the idealisation of $A$ by its canonical module. It is standard graded: every face $x_Sy_R$ of the whiskered complex extends to a facet $x_Sy_{[n]\setminus S}$ of size $n$, and for a monomial $m$ of degree below $n$ contained in a facet monomial $mm'$, the element $m'\cdot(mm')\dual$ equals $m\dual$, because $mm'$ is the only basis monomial of the form $m'w$. Hence $T$ is generated by $A_1$ and $(A_n)\dual$. Holleben and Nicklasson attach to a pure complex with facets $F_1,\dots,F_r$ the \emph{simplicial form} $F=\sum_kX_ku_{F_k}$ in new variables $X_k$, and study the Gorenstein algebra with Macaulay dual generator $F$. When $F$ is a \emph{Perazzo form} (the monomials $u_{F_k}$ are algebraically dependent, which holds when multiplication by a power of $L$ from degree one to the top degree is not surjective), this algebra is the Nagata idealisation $T$ in characteristic zero \cite[Lemmas~4.3 and~4.5]{HN}. For whiskered graphs it is moreover \emph{G-quadratic}, that is, its defining ideal has a quadratic Gr\"obner basis \cite[Theorem~4.4]{HN}.

Let $\phi\in(A_n)\dual$ be the coefficient of $u_{[n]}$ in the basis of Theorem~\ref{split}; in the original basis it is the coefficient of $y_{[n]}$. Put $\zeta=(L,\phi)\in T_1$. For a graded module $M$ write $M^{*D}$ for its reflected dual, $(M^{*D})_i=(M_{D-i})\dual$, and put $M_S=x_S\B_{[n]\setminus S}\otimes\kk$.

\begin{theorem}\label{ideal}
Multiplication by $\zeta$ on $T$ decomposes as
\begin{equation}\label{idealform}
 \B_{n+1}\ \oplus\bigoplus_{\varnothing\ne S\in\Delta}\left(M_S\oplus M_S^{*D}\right),
\end{equation}
where the first summand is the empty-face block together with its reflected dual, and each nonempty $M_S$ and $M_S^{*D}$ carries $U$ and its transpose. If $p=0$ or $p>D$, then $\zeta$ maximises the rank of every power $T_i\to T_j$ simultaneously; for $i+j\ge D$
\begin{equation}\label{coker}
 \dim\cok(\zeta^{j-i}:T_i\to T_j)=\dim\cok(L^{j-i}:A_i\to A_j),
\end{equation}
and the map $T_{D-j}\to T_{D-i}$ has the same rank as $T_i\to T_j$.
\end{theorem}
\begin{proof}
\emph{Splitting.} Multiplication by $\zeta$ sends $(b,\lambda)$ to $(Lb,\ L\lambda+b\phi)$. If $b\in M_S$ with $S\ne\varnothing$, every product $bw$ still contains $x_S$, so $(b\phi)(w)=\phi(bw)=0$; on these summands $\zeta$ acts by $U$ on $M_S$ and by the transpose of $U$ on $M_S^{*D}$, and the dual of $A$ splits accordingly. On $M_\varnothing=\B_n$ the pairing $(b,c)\mapsto\phi(bc)$ is perfect between complementary degrees, and $b\mapsto b\phi$ maps $\B_n$ isomorphically onto $M_\varnothing^{*D}$, raising degree by one. Writing $v\,b$ for $b\phi$ identifies $M_\varnothing\oplus M_\varnothing^{*D}$ with $\B_n\oplus v\B_n=\B_n[v]/(v^2)=\B_{n+1}$, on which $\zeta$ acts as $U+v$, because $L\cdot(b\phi)=(Lb)\phi$.

\emph{Cokernels.} Let $i+j\ge D$. The block $\B_{n+1}$ has twice-centre $D$ and is surjective. A reflected block $M_S^{*D}$ with $|S|=s\ge1$ occupies degrees $1,\dots,D-s$ and has twice-centre $n+2-s\le D$, so it is surjective too. Hence all cokernel of $\zeta^{j-i}$ comes from the blocks $M_S$, $S\ne\varnothing$, which are exactly the blocks of $A$ other than the empty-face block; that block is surjective on $A$ because $i+j\ge n+1$. This proves \eqref{coker}. For any $\zeta'=(h,\psi)\in T_1$, the projection $T\to A$ is a ring homomorphism, so it maps $\cok(\zeta'^{\,j-i})$ onto $\cok(h^{j-i})$, whose dimension is at least that of $\cok(L^{j-i})$ by Lemma~\ref{generic}. So $\zeta$ has the largest possible rank whenever $i+j\ge D$. The perfect pairing makes multiplication by $\zeta$ self-adjoint between $T_i\to T_j$ and $T_{D-j}\to T_{D-i}$, which gives the remaining cases.
\end{proof}

\begin{corollary}\label{tg}
Over a field of characteristic $p\ge0$, $T$ has the WLP if and only if $a\le b_n$ and either $p=0$ or $p>\lceil(n+1)/2\rceil$. It has the SLP if and only if $a\le1$ and either $p=0$ or $p>n+1$.
\end{corollary}
\begin{proof}
\emph{Sufficiency, WLP.} Let $a\le b_n$ and $p=0$ or $p>\lceil D/2\rceil$. Every block in \eqref{idealform} has size at most $D$, so by \eqref{wil} every adjacent block map has maximal rank. If $a\le1$ all blocks have twice-centre $D$. If $n$ is odd and $a=2$, the twice-centres are $D-1$, $D$ and $D+1$ with $D$ even; an odd sum $2i+1$ never lies strictly between $D-1$ and $D+1$, so all blocks are injective together or surjective together.

\emph{Sufficiency, SLP.} If $a\le1$ all blocks have twice-centre $D$, and for $p=0$ or $p>D$ every power of every block has maximal rank by \eqref{wil}.

\emph{Necessity, independence number.} We follow the projection argument of \cite[Theorem~4.6]{HN}. Suppose $a>b_n$, and let $i=\lceil n/2\rceil$, $j=i+1$. Then $i+j$ lies in the strip \eqref{strip}, with $s=b_n+1$ in Step~1 of Theorem~\ref{base}, so $h:A_i\to A_{i+1}$ has a cokernel for every linear form $h$, and so does every $\zeta'=(h,\psi)$ on $T$ after projection. If $T$ had the WLP, its Hilbert function would be unimodal (the first map that is not injective is surjective, and surjectivity then propagates upwards since $T$ is generated in degree one), and by symmetry about $D/2$ the map $T_i\to T_{i+1}$ would have to be surjective, because $i\ge D/2$ or $\dim T_i=\dim T_{i+1}$. This is a contradiction. For the SLP with $a\ge2$, the map $T_1\to T_n$ would have to be bijective since $\dim T_1=\dim T_n$, but its projection to $A$ misses the nonzero degree-$n$ part of a block $\B_{n-2}\langle2\rangle$, whose degree-one part is zero.

\emph{Necessity, characteristic.} For $\zeta'=(h,\psi)$ one has $\zeta'^{\,k}=(h^k,kh^{k-1}\psi)$, so $\zeta'^{\,p}=0$ in characteristic $p>0$. The SLP requires $\zeta'^{\,D}\ne0$, so $p>D$. The WLP together with symmetry and unimodality requires the maps $T_0\to T_1\to\cdots\to T_{\lceil D/2\rceil}$ to be injective, so $\zeta'^{\,\lceil D/2\rceil}\ne0$ and $p>\lceil D/2\rceil$.
\end{proof}

\begin{remark}
Let $G$ be non-complete, so $\alpha(G)\ge2$. The proof of Corollary~\ref{tg} shows that $L^{n-1}:(A_G)_1\to(A_G)_n$ is not surjective, so $F(w(G))$ is a simplicial Perazzo form \cite[Lemma~4.3]{HN} and, in characteristic zero, its Gorenstein algebra is $T$ \cite[Lemma~4.5]{HN}; it is G-quadratic by \cite[Theorem~4.4]{HN}. Corollary~\ref{tg} therefore says that this G-quadratic Perazzo algebra has the WLP in characteristic zero exactly when $n$ is odd and $\alpha(G)=2$. This extends \cite[Theorem~4.6 and Corollary~4.7]{HN}, which give failure when $\alpha(G)\ge n/3+2$.
\end{remark}

\section{Higher truncation exponents}\label{sec:higher}

This section shows that the change of variables also splits the algebras with higher powers of the variables considered by Holleben and Nicklasson \cite[Section~3]{HN}.

Fix integers $d_i\ge2$ and replace $x_i^2,y_i^2$ by $x_i^{d_i},y_i^{d_i}$ in the definition of $A_\Delta$, keeping $x_iy_i$ and $x_F$. Write $E=\sum_i(d_i-1)$ for the top degree. Holleben and Nicklasson proved the SLP for complete graphs when $p=0$ or $p>E$ \cite[Proposition~3.2]{HN} and the WLP for $a\le2$ when $E$ is odd and $p=0$ \cite[Proposition~3.5]{HN}.

\begin{theorem}\label{higher}
Over $\mathbb Z$ the operator $L$ splits into one shifted tensor product per face:
\[
 \bigoplus_{S\in\Delta}\Big(\bigotimes_{i=1}^n\mathbb Z[u_i]/(u_i^{\,d_i-\mathbf1_{i\in S}})\Big)\langle|S|\rangle ,
\]
with $L$ acting as $\sum_iu_i$. Let $h_S(j)$ be the dimension in degree $j$ of the summand for $S$. In characteristic zero the rank of $L^{j-i}$ from degree $i$ to degree $j$ is $\sum_S\min\{h_S(i),h_S(j)\}$. If all $d_i$ equal a common value $d$, then in characteristic zero that map fails to have maximal rank exactly when $E<i+j<E+a$; in particular the WLP holds exactly when $a\le b_E$ and the SLP exactly when $a\le1$.
\end{theorem}
\begin{proof}
\emph{Splitting.} Because $x_iy_i=0$, for $r\ge1$ we have $u_i^r=x_i^r+y_i^r$ and $x_iu_i^{r-1}=x_i^r$. So one pair of variables has the integral basis $1,u_i,\dots,u_i^{d_i-1},x_i,x_iu_i,\dots,x_iu_i^{d_i-2}$, unitriangular with respect to $1,x_i,\dots,x_i^{d_i-1},y_i,\dots,y_i^{d_i-1}$. The ideal generated by the $x_F$, $F\notin\Delta$, is spanned by the basis monomials whose $x$-support contains a nonface, since multiplying by a variable never removes an index from the $x$-support ($x_i\cdot x_i=x_iu_i$). Multiplication by $\sum u_i$ preserves the $x$-support $S$ and acts on the span of the $x_Su^{\beta}$ as on the stated tensor product.

\emph{Ranks.} Each summand is a monomial complete intersection, on which, in characteristic zero, $\sum_iu_i$ is a strong Lefschetz element \cite{Stanley}: its power from degree $i$ to degree $j$ is injective if $i+j$ is at most the twice-centre and surjective otherwise (as used in \cite[Proposition~3.2]{HN}). With Lemma~\ref{generic} this gives the rank formula. For unequal exponents only the rank formula is claimed: a summand can then have a plateau away from its centre, and the clean strip description can fail.

\emph{The strip for equal exponents.} The summand for $S$ has support $[|S|,E]$ and twice-centre $E+|S|$. When $n\ge2$ its Hilbert function is strictly increasing up to the middle (the two middle values may agree). To see this, list the chain lengths $e_i=d-1-\mathbf1_{i\in S}$ in decreasing order. The first two differ by at most one, so the product $(1+\cdots+t^{e_1})(1+\cdots+t^{e_2})$ has the property. If a symmetric polynomial $\sum_jh_jt^j$ of degree $R$ with the property is multiplied by $1+t+\cdots+t^{e}$ with $e\le R$, the new coefficients $g_j$ satisfy $g_j-g_{j-1}=h_j-h_{j-e-1}$. For $1\le j\le\lfloor(R+e)/2\rfloor$ we have $j\le R$, and either $j-e-1<0$ or $j-e-1$ is farther from $R/2$ than $j$ is; in both cases the difference is positive. Step~1 of Theorem~\ref{base} now applies verbatim with $E$ in place of $n$, using the face of size $i+j-E+1$. If $n=1$ then $a\le1$ and the strip is empty.
\end{proof}

We have not determined the prime-characteristic thresholds for Theorem~\ref{higher}; by the splitting they reduce to the known positive-characteristic behaviour of monomial complete intersections.

\section{Chordal graphs, paths and lollipops}\label{sec:chordal}

This section shows that removing a simplicial vertex splits the algebra $R(H)$ of an arbitrary graph, deduces a Lefschetz criterion for all chordal graphs, and applies it to paths and lollipop graphs.

Let $H$ be a finite simple graph, $R(H)$ as in the introduction (defined over $\mathbb Z$), and $L_H=\sum_vz_v$. Write $N[v]$ for the closed neighbourhood of $v$, and $H-X$ for the subgraph induced on the vertices outside $X$. A vertex $w$ is \emph{simplicial} if $N[w]$ is a clique. A graph is \emph{chordal} if every cycle of length at least four has a chord; by a theorem of Dirac, every chordal graph has a simplicial vertex (see \cite[p.~851]{FG}), and induced subgraphs of chordal graphs are chordal.

\begin{theorem}\label{chordal}
Let $w$ be a simplicial vertex of $H$ and $C=N[w]$. There is an integral decomposition of graded operators
\[
 R(H)\simeq\big(\B_1\otimes R(H-C)\big)\ \oplus\bigoplus_{v\in C\setminus\{w\}}R(H-N[v])\langle1\rangle ,
\]
where the operator on the tensor product is the sum of the two operators. Consequently, if $H$ is chordal, $R(H)$ splits into shifted Boolean blocks $\B_m\langle s\rangle$, one for each maximal independent set of $H$.
\end{theorem}
\begin{proof}
Replace $z_w$ by $u=\sum_{v\in C}z_v$ and keep the other variables. Since $C$ is a clique, all products of two variables from $C$, squares included, are relations; they span the same ideal in the new variables, giving $u^2=0$, $uz_v=0$ and $z_vz_{v'}=0$ for $v,v'\in C\setminus\{w\}$. A vertex outside $C$ is not adjacent to $w$, so all relations involving it are unchanged and $u$ has no relation with it. The quotient is again monomial. Its monomials containing no $z_v$ with $v\in C\setminus\{w\}$ are $u^\varepsilon z_I$ with $\varepsilon\in\{0,1\}$ and $I$ independent in $H-C$; those containing such a $z_v$ contain exactly one and have the form $z_vz_I$ with $I$ independent in $H-N[v]$ (note $C\subseteq N[v]$). Finally $L_H=u+\sum_{x\notin C}z_x$, which preserves both kinds of monomial and acts as stated.

For chordal $H$, the graphs $H-C$ and $H-N[v]$ are chordal, so the decomposition can be iterated down to the empty graph, whose algebra is $\B_0$; this produces Boolean blocks. A maximal independent set of $H$ contains exactly one vertex of the clique $C$ (if it contained none, $w$ could be added), and removing that vertex leaves a maximal independent set of $H-C$ or $H-N[v]$ respectively; this bijection gives the indexing by induction.
\end{proof}

Gasanova, Lundqvist and Nicklasson split the monomial basis of a monomial algebra along a colon ideal and obtain sufficient conditions for maximal rank from the two pieces \cite[Theorem~1]{GLN}; this is the tool behind \cite[Proposition~3.5]{HN}. Theorem~\ref{chordal} instead splits the operator itself, so maximal rank of $R(H)$ is equivalent to a condition on the summands. For a flag complex $\Delta=\operatorname{Ind}(G)$, the leaf $y_i$ of $w(G)$ is simplicial with $N[y_i]=\{x_i,y_i\}$, the substitution is $u_i=x_i+y_i$, and eliminating the $n$ leaves one after another reproduces Theorem~\ref{split}; there the remainders need not be chordal.

\paragraph{The block profile.}
Record the block $\B_m\langle s\rangle$ by the monomial $X^mY^s$ and let $P_H$ be the sum over all blocks (the \emph{block profile}). The proof of Theorem~\ref{chordal} gives, for a simplicial vertex $w$,
\begin{equation}\label{recursion}
 P_H=X\,P_{H-N[w]}+Y\!\!\sum_{v\in N[w]\setminus\{w\}}\!\!P_{H-N[v]},\qquad P_\varnothing=1 .
\end{equation}

\begin{proposition}\label{invariant}
Let $H$ be chordal.
\begin{enumerate}
\item The profile $P_H$ does not depend on the choice of simplicial vertices.
\item $P_{H\sqcup H'}=P_HP_{H'}$ for a disjoint union of chordal graphs.
\item The largest $m$ with $X^mY^s$ in $P_H$ equals $\alpha(H)$.
\end{enumerate}
\end{proposition}
\begin{proof}
(1) Over $\mathbb Q$, by \S\ref{sec:prelim} a block $\B_m\langle s\rangle$ contributes $\binom mh-\binom m{h-1}$ strings in degrees $s+h,\dots,s+m-h$ to the Jordan type of $L_H$. Every string with $h\ge1$ is strictly shorter than the string with $h=0$ of the same block. So a longest string $[s,s+m]$ of the Jordan type is the $h=0$ string of a block $\B_m\langle s\rangle$; removing the strings of that block and repeating recovers all blocks. The Jordan type of $L_H$ is determined by $H$, so the multiset of blocks is too.
(2) $R(H\sqcup H')=R(H)\otimes R(H')$ with operator $L_H\otimes1+1\otimes L_{H'}$, and $\B_m\langle s\rangle\otimes\B_{m'}\langle s'\rangle=\B_{m+m'}\langle s+s'\rangle$ with the summed operator. So $P_HP_{H'}$ is the profile of a block decomposition of $R(H\sqcup H')$, and by (1) it is the profile.
(3) In characteristic zero $L_H^k=k!\sum z_I$ over independent sets with $|I|=k$, so the largest $k$ with $L_H^k\ne0$ is $\alpha(H)$; on the other hand it is the largest $m$ of a block, since $U^m\ne0=U^{m+1}$ on $\B_m$.
\end{proof}

\begin{theorem}\label{chordalcrit}
Let $H$ be chordal with block profile $P_H$, and let $\mathcal T_H$ be the set of pairs $(m,s)$ occurring in $P_H$ (the \emph{block types}). For $i<j$, say that a block $(m,s)$ \emph{blocks injectivity} at $(i,j)$ if $s\le i\le s+m$ and $i+j>2s+m$, and \emph{blocks surjectivity} if $s\le j\le s+m$ and $i+j<2s+m$.
\begin{enumerate}
\item In characteristic zero, $L_H^{j-i}:R(H)_i\to R(H)_j$ fails to have maximal rank if and only if some block in $\mathcal T_H$ blocks injectivity and some block in $\mathcal T_H$ blocks surjectivity at $(i,j)$. Hence the WLP and SLP of $R(H)$ in characteristic zero depend only on $\mathcal T_H$.
\item In characteristic $p>0$, $R(H)$ has the WLP if and only if it has the WLP in characteristic zero and $p>\lceil\alpha(H)/2\rceil$, and it has the SLP if and only if it has the SLP in characteristic zero and $p>\alpha(H)$.
\end{enumerate}
\end{theorem}
\begin{proof}
By Lemma~\ref{generic} it suffices to study $L_H$. A block map is not injective exactly when its source is nonzero and larger than its target, which in characteristic zero means $s\le i\le s+m$ and $(i-s)+(j-s)>m$; symmetrically for surjectivity. A direct sum of maps of maximal rank has maximal rank if and only if the summands are not split between ``not injective'' and ``not surjective''. This gives (1).

For (2), every block has $m\le\alpha(H)$ by Proposition~\ref{invariant}(3). If $p>\lceil\alpha(H)/2\rceil$ the adjacent entries of \eqref{wil} are prime to $p$, so the ranks of $L_H$ between adjacent degrees are those of characteristic zero. If $2\le p\le\lceil\alpha(H)/2\rceil$, take a block with $m=\alpha(H)$ and degrees $s+p-1\to s+p$: as in Step~3 of Theorem~\ref{base} this block map lacks maximal rank, so it has both a kernel and a cokernel, and so does $L_H$. For the SLP, $p>\alpha(H)$ makes every entry of \eqref{wil} prime to $p$, and if $p\le\alpha(H)$ then $h^p=0$ for every linear form while $R(H)_p\ne0$.
\end{proof}

\paragraph{Paths.}
Let $P_n$ be the path with vertices $1,\dots,n$. Using the end vertex $1$ in \eqref{recursion} gives $P_{P_n}=XP_{P_{n-2}}+YP_{P_{n-3}}$ for $n\ge3$, with $P_{P_0}=1$, $P_{P_1}=X$, $P_{P_2}=X+Y$, and
\[
 \sum_{n\ge0}P_{P_n}t^n=\frac{1+Xt+Yt^2}{1-Xt^2-Yt^3}.
\]
For instance $P_{P_5}=X^3+2XY+Y^2$, of total dimension $8+2\cdot2+1=13$, the number of independent sets of $P_5$.

\begin{lemma}\label{pathtypes}
The block types of $P_n$ are exactly the pairs $(m,s)$ of nonnegative integers with $2m+3s\in\{n,n+1\}$.
\end{lemma}
\begin{proof}
This holds for $n\le2$ by the formulas above. For $n\ge3$, the types of $P_n$ are the types of $P_{n-2}$ with $m$ raised by one and those of $P_{n-3}$ with $s$ raised by one, which by induction are the pairs with $2m+3s\in\{n,n+1\}$ and $m\ge1$, or $s\ge1$, respectively; every pair other than $(0,0)$ has $m\ge1$ or $s\ge1$, and $(0,0)$ does not occur for $n\ge2$.
\end{proof}

\begin{corollary}\label{paths}
Over a field of characteristic $p$, $R(P_n)$ has the WLP if and only if $n\in\{1,2,\dots,7,9,10,13\}$ and either $p=0$ or $p>\lceil\lceil n/2\rceil/2\rceil$. It has the SLP if and only if $n\in\{1,2,\dots,7,9\}$ and either $p=0$ or $p>\lceil n/2\rceil$.
\end{corollary}
\begin{proof}
By Theorem~\ref{chordalcrit}(2) and $\alpha(P_n)=\lceil n/2\rceil$ it suffices to treat characteristic zero.

\emph{Failure.} Let $m_0=\lceil n/2\rceil$. The block $(m_0,0)$ occurs, and blocks injectivity at $(i,i+1)$ for $\lceil m_0/2\rceil\le i\le m_0$. A block $(m',s')$ with $s'\ge1$ blocks surjectivity at $(s'-1,s')$. By Lemma~\ref{pathtypes} such a block exists for every $s'$ with $3s'\le n+1$. So the WLP fails whenever some integer $s'$ satisfies $\lceil m_0/2\rceil+1\le s'\le\lfloor(n+1)/3\rfloor$, that is, whenever
\begin{equation}\label{pathwitness}
 \lceil n/4\rceil+1\le\lfloor(n+1)/3\rfloor .
\end{equation}
Since $\lceil n/4\rceil+1\le(n+7)/4$ and $\lfloor(n+1)/3\rfloor\ge(n-1)/3$, inequality \eqref{pathwitness} holds for all $n\ge25$, and direct inspection shows that for $n\le24$ it holds exactly when $n\in\{8,11,12\}$ or $n\ge14$.

\emph{The remaining cases.} For $n\in\{1,\dots,7,9,10,13\}$ the set of block types is given by Lemma~\ref{pathtypes}, and a finite check of the criterion of Theorem~\ref{chordalcrit}(1) over all adjacent pairs shows that no conflict occurs. Since the SLP implies the WLP, only these ten values need to be examined for the SLP, and the same finite check over all pairs $i<j$ shows that it holds for $n\le7$ and $n=9$ and fails for $n=10$ and $n=13$.
\end{proof}

The characteristic-zero WLP statement is the theorem of Nguyen and Tran \cite[Theorem~1.2(1)]{NT}, proved there by an induction on the independence polynomial supported by Macaulay2 computations.

\paragraph{Lollipops.}
The lollipop graph $L_{m,n}$ joins a complete graph $K_m$ to a path with $n$ vertices by a bridge from a vertex of $K_m$ to an end of the path \cite{TPS}; $L_{1,n}=P_{n+1}$. For $m\ge2$, a vertex of $K_m$ away from the bridge is simplicial, and \eqref{recursion} gives
\[
 P_{L_{m,n}}=X\,P_{P_n}+Y\big((m-2)P_{P_n}+P_{P_{n-1}}\big).
\]
So the block types of $L_{m,n}$ are the types of $P_n$ with $m$ raised by one, together with the types of $P_{n-1}$ (and, if $m\ge3$, of $P_n$) with $s$ raised by one. By Theorem~\ref{chordalcrit}(1) the characteristic-zero Lefschetz properties of $L_{m,n}$ are therefore the same for all $m\ge3$, which explains the uniform answer for $m\ge3$ in \cite{TPS}. Also $\alpha(L_{m,n})=\lceil n/2\rceil+1$ for $m\ge2$.

\begin{corollary}\label{lollipops}
Let $m\ge2$, $n\ge1$ and $\alpha=\lceil n/2\rceil+1$. Over a field of characteristic $p$, $R(L_{m,n})$ has the WLP if and only if either $p=0$ or $p>\lceil\alpha/2\rceil$, and $n\in\{1,\dots,5,7,8,11\}$ when $m=2$, or $n\in\{1,3,4,7\}$ when $m\ge3$. It has the SLP if and only if either $p=0$ or $p>\alpha$, and $n\in\{1,\dots,5,7\}$ when $m=2$, or $n\in\{1,3\}$ when $m\ge3$.
\end{corollary}
\begin{proof}
By Theorem~\ref{chordalcrit}(2) it suffices to treat characteristic zero. The block $(\lceil n/2\rceil+1,0)$ occurs and blocks injectivity at $(i,i+1)$ for $\lceil(\lceil n/2\rceil+1)/2\rceil\le i\le\lceil n/2\rceil+1$. For each $s'\ge0$ with $3s'\le n$, Lemma~\ref{pathtypes} gives a block $(m',s')$ of $P_{n-1}$, hence a block $(m',s'+1)$ of $L_{m,n}$, which blocks surjectivity at $(s',s'+1)$; for $m\ge3$ the same holds when $3s'\le n+1$, using $P_n$. So the WLP fails whenever
\[
 \big\lceil(\lceil n/2\rceil+1)/2\big\rceil\le\lfloor n/3\rfloor,\quad\text{or}\quad m\ge3\ \text{and}\ \big\lceil(\lceil n/2\rceil+1)/2\big\rceil\le\lfloor(n+1)/3\rfloor .
\]
The first inequality holds for all $n\ge23$, since its left side is at most $(n+5)/4$ and its right side at least $(n-2)/3$. For $n\le22$, direct inspection shows that these witnesses exist exactly when $n$ is not in the stated list, and a finite check of Theorem~\ref{chordalcrit}(1) shows that the listed cases have the WLP. For the SLP only the listed WLP cases need to be examined, and the finite check over all pairs $i<j$ gives the stated lists.
\end{proof}

The characteristic-zero WLP statement is \cite[Theorem~1.3]{TPS}, obtained there from a tensor-product criterion, snake-lemma arguments and induction. The method does not reach cycles $C_m$ with $m\ge4$ or tadpole graphs whose cycle has length at least four \cite{NT,TPS2}, since these are not chordal: eliminating simplicial vertices from such a tadpole graph stops at the cycle. (A tadpole graph with a triangle is the lollipop $L_{3,n}$.)

\section{Computations and scope}\label{sec:scope}

This section describes the independent computational checks and the limits of the results.

\paragraph{Checks in the original variables.}
The script \texttt{verify.py} forms multiplication matrices in the original monomial basis, without the change of variables, and computes ranks by exact elimination over prime fields. For every labelled graph on one to four vertices it compares the ranks of all powers of $L$ on $A_G$ and of $\zeta$ on $T$ with the predictions of Theorems~\ref{split} and~\ref{ideal} over $\mathbb F_2,\mathbb F_3,\mathbb F_5,\mathbb F_7$, and checks the WLP and SLP thresholds of Theorem~\ref{base} and Corollary~\ref{tg}: 8{,}440 rank comparisons on 75 graphs, all in agreement. A non-flag complex on four vertices (minimal nonfaces $\{1,2,3\}$ and $\{1,2,4\}$) adds 120 comparisons. For all 894 labelled chordal graphs on one to five vertices it compares 18{,}680 ranks with the block decomposition of Theorem~\ref{chordal}, and for higher exponents $(3)$, $(4,4)$, $(3,3,3)$, $(2,3,5)$, with and without a nonface, it makes 143 comparisons over $\mathbb F_{1009}$ with Theorem~\ref{higher}. Direct elimination for Example~\ref{ex8} over $\mathbb F_{1009}$ gives rank $386$ between spaces of dimensions $400$ and $406$.

\paragraph{Checks for Section~\ref{sec:chordal}.}
The script \texttt{chordal\_lefschetz.py} uses exact integer arithmetic only. It computes block profiles by \eqref{recursion}; compares the ranks they predict with 1{,}128 direct eliminations in the original variables for paths with at most twelve vertices and lollipops $L_{m,n}$ with $2\le m\le5$, $n\le7$; confirms on 300 random chordal graphs that two different elimination orders give the same profile and that profiles multiply over disjoint unions; evaluates the criterion of Theorem~\ref{chordalcrit} for paths with $n\le120$ and lollipops with $m\le8$, $n\le60$, in characteristic zero, and at every prime $p\le43$ for paths with $n\le40$ and every prime $p\le29$ for lollipops with $n\le25$; and checks the witness inequalities of Corollaries~\ref{paths} and~\ref{lollipops} for every $n\le5000$. Its output reproduces the published characteristic-zero lists of \cite{NT} and \cite{TPS} on these ranges and gives the SLP lists stated above. The finite checks invoked in the proofs of Corollaries~\ref{paths} and~\ref{lollipops} are the cases $n\le24$ and $n\le22$ of this computation.

\paragraph{Scope.}
The splitting needs a private leaf at every vertex (or, for $R(H)$, a simplicial vertex at each step). For whiskered complexes the results hold over every field. For chordal graphs the profile is computed by \eqref{recursion}, which may branch; for paths and lollipops it collapses to the explicit recursions above. Partially whiskered graphs, very well-covered graphs in general \cite[Example~3.18]{HN}, cycles and tadpoles, and the classification of Perazzo algebras with the WLP remain outside these methods.

\paragraph{AI use and data.}
This manuscript was prepared with the assistance of AI systems, which were used for the mathematical development, the verification scripts and the drafting; the proofs were checked line by line and the scripts were rerun from a clean copy. The scripts and their recorded outputs accompany the manuscript.

\begin{thebibliography}{99}
\bibitem{Cooper}
S. M. Cooper, S. Faridi, T. Holleben, L. Nicklasson and A. Van Tuyl (2024).
The weak Lefschetz property of whiskered graphs. In \emph{Lefschetz Properties}, Springer INdAM Series, 97--110.
\url{https://doi.org/10.1007/978-981-97-3886-1_5}
\bibitem{FG}
D. R. Fulkerson and O. A. Gross (1965). Incidence matrices and interval graphs. \emph{Pacific Journal of Mathematics} \textbf{15}(3), 835--855.
\url{https://doi.org/10.2140/pjm.1965.15.835}
\bibitem{GKMM}
E. Ghorbani, G. B. Khosrovshahi, Ch. Maysoori and M. Mohammad-Noori (2008). Inclusion matrices and chains. \emph{Journal of Combinatorial Theory, Series A} \textbf{115}(5), 878--887.
\url{https://doi.org/10.1016/j.jcta.2007.09.002}
\bibitem{GLN}
O. Gasanova, S. Lundqvist and L. Nicklasson (2022). On decomposing monomial algebras with the Lefschetz properties. \emph{Journal of Pure and Applied Algebra} \textbf{226}(6), 106968.
\url{https://doi.org/10.1016/j.jpaa.2021.106968}
\bibitem{HN}
T. Holleben and L. Nicklasson (2026). Roller Coaster Gorenstein algebras and Koszul algebras failing the weak Lefschetz property. \emph{Journal of Pure and Applied Algebra} \textbf{230}(4), 108238.
\url{https://doi.org/10.1016/j.jpaa.2026.108238}
\bibitem{NT}
H. D. Nguyen and Q. H. Tran (2024). The weak Lefschetz property of Artinian algebras associated to paths and cycles. \emph{Acta Mathematica Vietnamica} \textbf{49}(3), 523--544.
\url{https://doi.org/10.1007/s40306-024-00549-1}
\bibitem{Stanley}
R. P. Stanley (1980). Weyl groups, the hard Lefschetz theorem, and the Sperner property. \emph{SIAM Journal on Algebraic Discrete Methods} \textbf{1}(2), 168--184.
\url{https://doi.org/10.1137/0601021}
\bibitem{TPS}
Tran Quang Hoa, Nguyen Duy Phuoc and Tran Nguyen Thanh Son (2025). The weak Lefschetz property for tensor products of Artinian monomial algebras and its applications to lollipop graphs. arXiv:2510.21569.
\url{https://arxiv.org/abs/2510.21569}
\bibitem{TPS2}
Tran Quang Hoa, Nguyen Duy Phuoc and Tran Nguyen Thanh Son (2026). Independence polynomials and the weak Lefschetz property for tadpole graphs. arXiv:2609.13888.
\url{https://arxiv.org/abs/2609.13888}
\bibitem{Wilson}
R. M. Wilson (1990). A diagonal form for the incidence matrices of $t$-subsets vs.\ $k$-subsets. \emph{European Journal of Combinatorics} \textbf{11}(6), 609--615.
\url{https://doi.org/10.1016/S0195-6698(13)80046-7}
\bibitem{Zeng}
Z. Zeng (2026). The whiskered-graph weak Lefschetz conjecture for graphs on eight vertices. SSRN preprint.
\url{https://doi.org/10.2139/ssrn.7385138}
\end{thebibliography}
\end{document}
